% ===========================================================================
% Section: Related Work
% ===========================================================================
\section{Related Work}
\label{sec:related}

\subsection{3D Gaussian Splatting for SLAM}

Since the introduction of 3D Gaussian Splatting~\cite{kerbl2023_3dgs},
numerous works have integrated Gaussian representations into SLAM systems.
SplaTAM~\cite{keetha2024_splatam} and MonoGS~\cite{matsuki2024_monogs}
demonstrated monocular and RGB-D Gaussian SLAM, respectively.
Photo-SLAM~\cite{huang2024_photoslam} combined ORB features with a
Gaussian map for real-time operation.
\gaussianlic~\cite{lang2026_gaussianlic2} advanced the field by fusing
LiDAR, inertial, and camera data for robust pose estimation and
photo-realistic mapping, achieving real-time performance on an NVIDIA
RTX~3090/4090.

Several follow-up works have explored different aspects of Gaussian SLAM:
Scaffold-GS~\cite{lu2024_scaffoldgs} introduced anchor-based structures
for efficient rendering, LIV-GaussMap~\cite{hong2024_livgaussmap}
focused on LiDAR-inertial-visual fusion, and
GSLIVM~\cite{xie2024_gslivm} addressed large-scale indoor-outdoor scenes.
Recent surveys~\cite{yan2024_gsslam, chen2025_egsslam} provide
comprehensive reviews of the rapidly growing field.

A common limitation across these systems is their reliance on ROS for
inter-module communication, restricting deployment to Linux and hindering
integration with commercial sensor SDKs.

\subsection{Multi-Modal Sensor Fusion for Depth Estimation}

Multi-modal depth fusion has been extensively studied in the robotics and
computer vision communities.
LiDAR-camera fusion~\cite{zheng2025_wildgsslam} leverages the complementary
characteristics of LiDAR (accurate but sparse) and cameras (dense but
ambiguous in depth).
iToF sensors provide dense depth at moderate range but suffer from
multi-path interference and noise at depth discontinuities~\cite{rang2018_4dof}.
Stereo matching via Semi-Global Matching (SGM)~\cite{hirschmuller2008_sgm}
produces dense disparity maps but fails in textureless regions.

Confidence-based fusion strategies have been proposed to combine these
modalities: Ramezani \emph{et al.}~\cite{ramezani2022_wildcat} used a
sliding-window approach for LiDAR-inertial-camera fusion, while
Coco-LIC~\cite{lang2023_cocolic} employed continuous-time B-spline
representation for tightly-coupled multi-modal odometry.
Our work adopts a simpler confidence-weighted fusion strategy suitable
for offline batch processing, where the SLAM front-end has already
produced accurate poses.

\subsection{Depth Completion}

Depth completion---filling holes in sparse depth maps---has been
addressed by learning-based methods.
SPNet, used in \gaussianlic, takes RGB and sparse depth as input and
produces a dense depth map via a TensorRT-accelerated neural network.
Alternative approaches include NLSPN~\cite{park2020_nlspn} and
GuideNet~\cite{tang2020_guidenet}, which use learned spatial propagation.
In \winlic, we retain the SPNet-based completion for algorithm fidelity,
but the modular architecture allows easy substitution of alternative
completion networks.
